\documentclass[11pt]{article}
\usepackage[T1]{fontenc}
\usepackage[utf8]{inputenc}
\usepackage[margin=23mm]{geometry}
\usepackage{amsmath,amssymb,booktabs,array,longtable}
\usepackage{xurl}
\usepackage{hyperref}
\hypersetup{colorlinks=true,linkcolor=blue,citecolor=blue,urlcolor=blue}
\newcommand{\PE}{\mathrm{PE}}
\newcommand{\E}{\mathbb{E}}
\newcommand{\Var}{\operatorname{Var}}
\newcommand{\code}[1]{\path{#1}}
\providecommand{\SiPMGeneratedParameters}{}
\InputIfFileExists{latex_generated/sipm_response/parameters.tex}{}{}
\title{SiPM charge response and provisional ASIC readout\\for NAUSICAA electroluminescence}
\author{NAUSICAA model and implementation note}
\date{4 November 2026}
\begin{document}
\maketitle

\begin{abstract}
The response begins with incident VUV photon counts at each SiPM in
250~ns buckets. A stochastic microcell model samples photodetection and dark
pulses, follows the charge recovery of each microcell, and saves the sensor
charge before applying electronics. The provisional ASIC reads bucket-average
current, adds independent Gaussian noise of 0.5~photoelectron equivalents per
bucket, and clips the output above 1~mA. This note states the adopted inputs,
physical approximations and output contract, and derives the relation to
AXEL's subsequent correction for SiPM nonlinearity. Numerical validation and a
small pilot precede a full-sample response calculation.
\end{abstract}

\section{Input, terminology and provenance}
The input is the number $N_{\gamma,cb}$ of photons incident on SiPM channel
$c$ in bucket $b$ of width $\Delta t$. The existing four HDF5 inputs represent
\code{bb0nu} and \code{single_electrons}, each with pure-Xe and Xe--15\% He
transport and electroluminescence models. The gas files are alternative
realizations of the detector response for the same underlying event sample.
The current production uses $\Delta t=0.25~\mu\mathrm{s}=250$~ns, as recorded
in the embedded detector card. The input index and exact source paths are
specified in \code{docs/sipm_electronics_handoff.md}. Each source has its own
content identity, event identifiers and provenance.

The upstream calculation includes electron diffusion, attachment,
electroluminescence production and geometrical photon collection. Its adopted
EL channel is 5~mm in diameter and depth, at 15~bar and
$E/P=3~\mathrm{kV\,cm^{-1}\,bar^{-1}}$. The adopted uniform emission along the
electron path motivates a uniform conditional photon arrival time within each
bucket. This is a timing approximation at the response interface. It cannot
restore sub-bucket timing that was not saved by the producer. The response
preserves the producer's signed bucket index and geometrical reference.

A \emph{primary photodetection} is the successful conversion of an incident
photon into an avalanche trigger. A \emph{dark pulse} is an avalanche initiated
without an incident signal photon. The integer number of triggers differs
from charge expressed in \emph{photoelectron equivalents} (PE): one PE denotes
the calibrated charge $q_0$ of a fully recovered avalanche. A partially
recovered avalanche contributes a fraction of one PE. We use \emph{microcell}
for an avalanche element inside a SiPM and \emph{channel} for a complete SiPM.

\section{Sensor reference and adopted operating point}
Hamamatsu's AXEL presentation describes a custom VUV-sensitive device based on
the S13370 family, with 50~$\mu$m microcell pitch, TSV connections and a
64-channel assembly on a flexible polyimide substrate with a PEEK
stiffener~\cite{Ghassemi2024MPPC}. The displayed PDE curve was taken at
$-25\,^{\circ}\mathrm{C}$ and 4~V overvoltage; it is not a calibration of our
room-temperature design. The AXEL prototype analysis names the standard
S13370-3050CN~\cite{Akiyama:2025dia}. That standard device has a
$3\times3~\mathrm{mm^2}$ active area and 3600 microcells. Our 5~mm diameter
active disk follows the optical geometry already used in photon production.
Its effective microcell count is a geometrical design approximation.

\begin{table}[htbp]
\centering\small
\begin{tabular}{p{0.30\linewidth}p{0.23\linewidth}p{0.37\linewidth}}
\toprule
Parameter & Adopted input & Status or origin\\
\midrule
Bucket width $\Delta t$ & 250 ns & User design input; detector card\\
Photon detection efficiency $p$ & 0.30 & User design input\\
Active disk diameter $d$ & 5 mm & Existing optical geometry\\
Microcell pitch $a$ & 50 $\mu$m & AXEL device family reference\\
Microcell count $M$ & $\operatorname{round}[\pi(d/2)^2/a^2]$ & Effective uniform-cell geometry\\
Fully recovered gain $G$ & $2.55\times10^6$ & Standard-device reference\\
Recovery time $\tau$ & 73.4 ns & User adopts AXEL measured mean\\
Reference dark rate $R_0$ & $10^6~\mathrm{s^{-1}}$ per 9 mm$^2$ & Typical reference at $25\,^{\circ}$C, 4 V\\
ASIC noise $\sigma_{\PE}$ & 0.5 PE RMS/bucket & User provisional input\\
ASIC upper range $I_{\max}$ & 1 mA & User provisional input\\
\bottomrule
\end{tabular}
\caption{Baseline inputs. Runtime choices reside in versioned TOML cards.
Gain and dark-rate references are provisional standard-device values;
they do not characterize a fabricated 5~mm circular device.}
\label{tab:inputs}
\end{table}

The gain and dark-rate reference comes from the search-indexed manufacturer
specification KSX-I50133-E~\cite{Hamamatsu2017VUV4Indexed}. The original PDF
remains to be recovered; this source limitation is recorded in the
bibliography and specification note. The table gives a typical dark rate of
1~MHz and a maximum of 3~MHz for the standard device at $25\,^{\circ}$C and
4~V overvoltage. The baseline takes the typical value as an independent
primary-dark-avalanche rate. Its measurement definition must be checked
before introducing separately sampled correlated avalanches. The same source
lists 3\% crosstalk; the baseline probability is nevertheless set to zero as an
explicit modeling approximation. A suitable room-temperature afterpulse
probability and delay distribution has not been established.

The PDE includes optical fill factor and avalanche initiation probability in
its calibrated definition. The reference dark rate also already reflects the
reference device's geometry. Neither receives an additional fill-factor
multiplier. Temperature and overvoltage are recorded as reference operating
conditions; a transfer law for these parameters is not assumed.

\subsection{Parameters exported from the saved calculation}
The following values are generated from the pilot HDF5 output. In particular,
the dark-rate scaling, one-PE charge, noise-current conversion and equivalent
current range are evaluated by the Julia calculation using the selected cards.
\begin{center}\small
\SiPMGeneratedParameters
\end{center}

\section{Stochastic microcell response}
\subsection{Detection, arrival times and illumination}
At fixed incident count, sample the number of primary photodetections as
\begin{equation}
 N_{s,cb}\mid N_{\gamma,cb}\sim
 \operatorname{Binomial}(N_{\gamma,cb},p).
 \label{eq:detection}
\end{equation}
Conditional on $N_{s,cb}$, draw each time independently and uniformly over
$[t_b,t_b+\Delta t)$, and assign each detection independently to one of the
$M$ microcells with probability $1/M$. This is equivalent to sampling times
and microcell labels for every incident photon and retaining independent
Bernoulli detections, while avoiding work for undetected photons. Uniform
illumination is an additional sensor approximation: the saved input resolves
SiPM channels, not the illumination pattern inside one sensor.

The baseline holds $p$ fixed during microcell recovery. Consequently a
successful trigger at short delay still initiates an avalanche, whose charge
is reduced by the law below. A model in which the initiation probability also
depends on recovered overvoltage would need an additional calibrated function
$p(r)$; the AXEL mean correction alone does not determine that function.

\subsection{Recovery and avalanche charge}
Let $t_{\mathrm{last},j}$ be the time of the previous avalanche in microcell
$j$. Just before a trigger at $t_k$, its recovered fraction and charge are
\begin{align}
 r_j(t_k^-)&=1-\exp[-(t_k-t_{\mathrm{last},j})/\tau],\\
 q_k&=q_0 r_j(t_k^-),\qquad q_0=eG.
 \label{eq:recovery}
\end{align}
After the avalanche set $t_{\mathrm{last},j}=t_k$. An explicitly fully recovered
microcell has $t_{\mathrm{last},j}=-\infty$. The default dark-equilibrium
initial state is specified below. The state is carried continuously
between buckets; a bucket boundary does not recharge the sensor.
Equation~\eqref{eq:recovery} follows the exponential gain-recovery model used
by AXEL~\cite{Ban2020AXELPrototype}. Charge accumulation and exponentials use
SI units and Float64 arithmetic in the implementation.

The time constant is not a fixed dead interval. At a delay $\tau$, the
recovered charge is $(1-e^{-1})q_0$; after $2\tau$ and $3\tau$ the fractions
are $1-e^{-2}$ and $1-e^{-3}$. Full recovery is approached asymptotically.
The probability to initiate a second avalanche and its available charge are
distinct quantities; this baseline specifies the former by a constant PDE
and the latter by equation~\eqref{eq:recovery}.

For all avalanches assigned to channel $c$ and bucket $b$,
\begin{align}
 Q_{s,cb}&=\sum_{k\in\mathrm{signal}}q_k,&
 Q_{d,cb}&=\sum_{k\in\mathrm{dark}}q_k,\\
 Q_{cb}&=Q_{s,cb}+Q_{d,cb},&
 m_{cb}&=Q_{cb}/q_0,&
 \overline I_{cb}&=Q_{cb}/\Delta t.
 \label{eq:bucket-charge}
\end{align}
The signal and dark sums label the cause of each avalanche. Because both
populations change the same microcell state, $Q_s$ is affected by preceding
dark avalanches and $Q_d$ by preceding signal avalanches. These saved
components provide charge accounting, rather than independent responses that
could be added after simulating two isolated sensors.

\subsection{Where the fluctuations enter}
For known microcell labels and arrival times, equation~\eqref{eq:recovery}
is deterministic. In an ensemble, the labels, times, primary photodetections
and dark counts vary. Two detected photons may strike different cells and
both produce full charge, or strike the same cell and produce
$q_0+q_0[1-\exp(-\delta t/\tau)]$ if it was initially recovered. Their charge
therefore fluctuates even at fixed detected-photon count. Additional
intrinsic single-avalanche gain dispersion is set to zero in the baseline.

At low occupancy and in the absence of dark pulses, the recovered-cell limit
provides a useful check:
\begin{equation}
 \E[Q\mid N_\gamma]=q_0pN_\gamma,\qquad
 \Var(Q\mid N_\gamma)=q_0^2p(1-p)N_\gamma.
 \label{eq:low-occupancy}
\end{equation}
This is the variance conditional on the supplied photon count. Upstream
photon-production and collection fluctuations are already represented by the
input event ensemble. They must not be resampled as an extra Poisson input.

\section{Dark pulses and the sensor's dark current}
The adopted area scaling at fixed temperature, overvoltage and technology is
\begin{equation}
 A=\pi(d/2)^2,\qquad
 R_d=R_0\frac{A}{A_0},\qquad A_0=9~\mathrm{mm^2}.
 \label{eq:dark-area}
\end{equation}
For each selected channel and acquisition bucket draw
$N_{d,cb}\sim\operatorname{Poisson}(R_d\Delta t)$, with uniform times and
uniform microcell labels. Merge these triggers with photodetections in time
order before updating the cell states. Equivalently one can generate a
homogeneous Poisson process over the acquisition interval. The time-window
and initial-state choices are part of the response card and provenance.
Dark pulses occur in buckets with zero incident photons as well as signal
buckets.

The default acquisition uses every channel illuminated anywhere in the event.
All selected channels are simulated over the same interval: the global first
to last occupied photon bucket, with four additional buckets on each side
(1~$\mu$s at the default sampling). This defines the channel set for the
pilot; it does not include the unilluminated remainder of the detector.
The sensor card can add explicitly listed channels. For an empty photon
history it defines the first and last acquisition bucket, supporting a
dark-only acquisition. No threshold is applied to the output samples.

At acquisition start $t_0$, initialize each microcell with an independent age
$U_j\sim\operatorname{Exponential}(\text{mean}=M/R_d)$ and set
$t_{\mathrm{last},j}=t_0-U_j$. This is the stationary backward-recurrence
age for the assumed Poisson dark triggers. It represents the dark activity
before the recorded window without storing those earlier pulses. At zero
dark rate the cells start fully recovered. This initialization assumes the
signal begins within the recorded photon interval, with only dark activity
preceding it.

In the dilute, recovered-cell limit, a dark-only exposure of duration $T$
has
\begin{equation}
 \E[Q_d]=q_0R_dT,\qquad \Var(Q_d)=q_0^2R_dT,
 \qquad \E[I_d]=q_0R_d.
 \label{eq:dark-current}
\end{equation}
Thus dark current is the time-averaged charge carried by random dark pulses,
and the pulse fluctuations survive subtraction of its expectation. At finite
occupancy recovery reduces both pulse amplitudes and the average dark
current; the stationary result is derived below. A continuous leakage current
would be a separate baseline term, currently taken as zero. Independent
dark-count processes across channels and times are an approximation to the
primary thermal avalanche population. Optical crosstalk and afterpulsing
would add correlated triggers and require separate calibrated parameters.

\section{Current readout and the provisional ASIC}
The SiPM supplies avalanche current pulses; a voltage signal arises only
through the attached electrical network. For the current bucket model, each
avalanche contributes its integrated charge to its bucket, so the observable
is $\overline I=Q/\Delta t$. This convention preserves
\begin{equation}
 \sum_b \overline I_{cb}\Delta t=\sum_k q_k
 \label{eq:charge-preservation}
\end{equation}
for avalanches inside the acquisition interval. The model specifies neither
an instantaneous peak current nor a shaped voltage waveform. A finite pulse
shape and front-end transfer function would redistribute charge across
buckets and require a revised readout calculation with explicit boundary
losses.

With known zero baseline, the adopted ASIC model is
\begin{align}
 \epsilon_{cb}&\sim\mathcal N(0,\sigma_{\PE}^2),\\
 I^{\mathrm{raw}}_{cb}&=\overline I_{cb}
       +\epsilon_{cb}\frac{q_0}{\Delta t},\\
 I^{\mathrm{out}}_{cb}&=\min(I^{\mathrm{raw}}_{cb},I_{\max}),\\
 s_{cb}&=\boldsymbol{1}\{I^{\mathrm{raw}}_{cb}>I_{\max}\}.
 \label{eq:asic}
\end{align}
Noise is independent between channels and buckets. It is added before the
upper range limit. Negative baseline-subtracted samples are retained; the
model has no lower clipping threshold. Save the sampled noise, current before
clipping, output current and saturation flag. The independent sensor output
remains available for changing or diagnosing the electronics model.

The conversion and range in PE units follow from the selected gain:
\begin{equation}
 I_{1\PE}=\frac{q_0}{\Delta t},\quad
 \sigma_I=\sigma_{\PE}\frac{q_0}{\Delta t},\quad
 m_{\max}^{\mathrm{ASIC}}=\frac{I_{\max}\Delta t}{q_0}.
 \label{eq:range}
\end{equation}
The fixed hardware input is 1~mA. Its PE equivalent changes with $G$ and
$\Delta t$. This specification supersedes the earlier approximate 300-PE
range discussed before the ASIC current limit was adopted. It always applies
to the sensor's measured charge/current response, after photon detection,
with the Gaussian term added as in equation~\eqref{eq:asic}.

\subsection{Relation to AXEL readout and noise measurements}
AXEL's front-end electronics use a 5~MS/s physics readout, with a 17~ns
first-order stage followed by an approximately 400~ns second-order
Sallen--Key stage. A 40~MS/s auxiliary readout supports single-PE dark-pulse
gain calibration~\cite{Nakamura2020AxFEB}. Its physical sampling interval is
200~ns, whereas the adopted NAUSICAA interval is 250~ns. These transfer
functions do not define our ASIC.

The 2025 AXEL paper reports a mean baseline standard deviation of 0.46~ADC
counts during operation and 0.45 without high voltage, with a 12-bit ADC and
0.488~mV per ADC count~\cite{Akiyama:2025dia}. The 2024 analysis uses a
3.5-ADC threshold, corresponding to 3.8 detected-photon units, and extends
integration by five samples on each side of the threshold crossings. It
also considers a baseline-offset uncertainty within one ADC count
separately~\cite{Yoshida2024AXELResponse}. A persistent baseline offset and
independent sample noise propagate differently when integrating an event.

Combining a PE calibration from one run with baseline RMS from another only
supplies a rough comparison. Our 0.5-PE Gaussian term is the adopted ASIC
input, additional to explicitly generated sensor dark pulses. It is not a
measured separation of AXEL's baseline into dark, analog and digitization
components. Shaping-induced correlations, offset drift and ADC quantization
remain extensions requiring ASIC information.

\section{Mean linearity and subsequent signal recovery}
\subsection{AXEL's stationary mean response}
The AXEL correction can be derived from the same exponential charge law
under steady uniform illumination~\cite{Ban2020AXELPrototype}. Let the
combined successful signal and dark triggers in each microcell form a
stationary Poisson process of rate $\lambda$. Its interval density is
$f(u)=\lambda e^{-\lambda u}$ for $u\geq0$. Averaging the recovered fraction
at avalanche times gives
\begin{equation}
 \E[r]=\int_0^\infty(1-e^{-u/\tau})\lambda e^{-\lambda u}\,du
       =\frac{1}{1+\lambda\tau}.
 \label{eq:mean-recovery}
\end{equation}
For expected primary avalanche count $n=M\lambda\Delta t$, the mean charge
in PE units is
\begin{equation}
 m=\frac{n}{1+kn},\qquad k=\frac{\tau}{M\Delta t}.
 \label{eq:forward-mean}
\end{equation}
The low-intensity limit is $m\simeq n-kn^2$. The large-rate stationary
current approaches $Mq_0/\tau$. This stationary asymptote is not a bound on
the first short bucket of an initially charged sensor, which also contains
stored charge. Sensor nonlinearity depends on occupancy per microcell and
recovery time; electronics clipping depends on the configured current limit.

Dark-only stationary current follows by setting $\lambda=R_d/M$:
\begin{equation}
 \E[I_d]=\frac{q_0R_d}{1+R_d\tau/M}.
 \label{eq:dark-stationary}
\end{equation}
Equations~\eqref{eq:mean-recovery}--\eqref{eq:dark-stationary} assume constant
rate, independent triggering and fixed initiation probability. The simulation
conditions on the supplied finite counts and follows state across time-varying
buckets, so the stationary expression is a validation limit and a candidate
reconstruction law, not an identity for every event.

\subsection{Compact recovery-state and charge moments}
\label{sec:compact-recovery}
The following local derivation establishes the basis for a faster stochastic
bucket response. It uses the implemented assumptions: uniform independent
microcell selection, fixed initiation probability, exponential recovery, and
reset to zero recovery after every successful trigger. The current production
implementation follows individual avalanches; the compact model below is a
proposed implementation strategy.

\paragraph{Exact closure of the mean state.}
Let $r_j\in[0,1]$ be the recovered fraction of microcell $j$ and define
$R=\sum_{j=1}^M r_j$. Between avalanches,
\begin{equation}
 \frac{dR}{dt}=\frac{M-R}{\tau}.
 \label{eq:total-recovery-drift}
\end{equation}
At an avalanche, a uniformly selected cell $J$ contributes $r_J$ PE and
resets to zero. Conditional on the entire state $\boldsymbol r$,
\begin{align}
 \E[\Delta m\mid\boldsymbol r]&=\frac{R}{M},\\
 \E[R^+\mid\boldsymbol r]&=\left(1-\frac1M\right)R^-.
 \label{eq:mean-state-jump}
\end{align}
Consequently, conditional on a sequence of arrival times, the ensemble mean
of $R$ and the expected collected charge evolve using only the mean initial
$R$. This closure is exact for these assumptions. Averaging over arrival
times preserves it. A distribution of recovery states is unnecessary for
calculating the mean charge.

\paragraph{The specific stochastic contribution.}
Random microcell selection contributes the conditional avalanche-charge
variance
\begin{equation}
 \Var(\Delta m\mid\boldsymbol r)
 =\frac1M\sum_{j=1}^M r_j^2-\left(\frac{R}{M}\right)^2.
 \label{eq:cell-choice-variance}
\end{equation}
A deterministic charge assignment $\Delta m=R/M$ removes this contribution.
An aggregate stochastic response can retain it through charge moments.
The same avalanche changes the outgoing recovery state, which creates a
correlation between collected charge and the response to subsequent light.

\paragraph{Closed moment equations in the Poisson limit.}
For independent Poisson triggers of deterministic rate $\lambda(t)$ per
microcell, let $u_1=\E[r]$ and $u_2=\E[r^2]$. The recovery drift and reset give
\begin{align}
 \dot u_1&=\frac{1-u_1}{\tau}-\lambda u_1,\\
 \dot u_2&=\frac{2(u_1-u_2)}{\tau}-\lambda u_2.
 \label{eq:recovery-moments}
\end{align}
At constant rate, their stationary values are
\begin{equation}
 u_1=\frac{1}{1+\lambda\tau},\qquad
 u_2=\frac{2}{(1+\lambda\tau)(2+\lambda\tau)}.
\end{equation}
The first reproduces equation~\eqref{eq:mean-recovery}.

Let $q$ denote accumulated charge in PE units from one microcell during the
current bucket, and write $a=\E[q]$, $b=\E[qr]$ and $c=\E[q^2]$.
At a trigger $(q,r)\mapsto(q+r,0)$, yielding
\begin{align}
 \dot a&=\lambda u_1,\\
 \dot b&=\frac{a-b}{\tau}-\lambda b,\\
 \dot c&=\lambda(2b+u_2).
 \label{eq:charge-state-moments}
\end{align}
Start each bucket with $a=b=c=0$ and the incoming recovery moments.
At its end, the single-cell charge variance is $c-a^2$ and the covariance
of collected charge with final recovery is $b-a u_1$. For independent
microcells, the total mean and variance are $Ma$ and $M(c-a^2)$, and the
covariance with total outgoing recovery is $M(b-a u_1)$.
Thus a finite set of moment equations determines these observables in the
Poisson model, including the effect of collection throughout the bucket.

\paragraph{Production interpretation.}
The implemented optical counts are binomial conditional on saved photons;
arrival times are uniform conditional on those counts. A Poisson replacement
would change the prescribed count fluctuations. The compact production model
should preserve those count draws and derive or calibrate its conditional
charge and outgoing-state moments accordingly. The equations above provide
an exact analytical limit and a practical starting point for that calculation.
Its stochastic sampler must reproduce the charge variance and its covariance
with outgoing recovery. Validation observables are bucket means and variances,
neighbouring-bucket covariance, and the total-event charge distribution.
A detailed recovery-state distribution should be introduced only if these
observables demonstrate a consequential failure of the compact model.

\subsection{Conditional-count statistical implementation}
\label{sec:statistical-implementation}
The prototype in \code{src/ion_signal/sipm_statistical.jl} implements a
count-preserving discretization of the moment model. It retains the binomial
photodetection and Poisson dark-count draws. The state carried between buckets
is $(R,S)$, with $S=\sum_j r_j^2$. Dark-equilibrium initialization samples the
cell recoveries once and reduces them to these two totals.

\paragraph{Exact conditional occupancy moments.}
For $k$ simultaneous triggers assigned independently to $M$ cells, write
\begin{equation}
 a=(1-1/M)^k,\qquad b=(1-2/M)^k,\qquad p=1-a.
\end{equation}
Each occupied cell contributes its available charge once. For interval charge
$Q$ in PE equivalents,
\begin{align}
 \E[Q\mid R,S,k]&=pR,\\
 \Var(Q\mid R,S,k)&=(a-b)S+(b-a^2)R^2,\\
 \E\!\left[\sum_{j\,\mathrm{hit}}r_j^2\mid R,S,k\right]&=pS.
 \label{eq:conditional-occupancy}
\end{align}
These expressions follow from the single-cell occupancy probability $p$ and
the two-cell occupancy covariance $b-a^2$. They condition on the actual
trigger count. The implementation uses a numerically stable equivalent of
the variance expression and handles zero and one trigger explicitly.

\paragraph{Charge and outgoing state.}
The interval charge uses a bounded beta draw matching the first two moments
in equation~\eqref{eq:conditional-occupancy}. Its support respects the total
available charge, trigger count and the maximum cell recovery allowed by
$(R,S)$. Zero variance gives a deterministic draw; maximal bounded variance
uses a two-point distribution. The outgoing first moment is $R^+=R^--Q$,
so output charge and recovery loss are correlated by construction.
For a single trigger the second-moment loss is exactly $Q^2$. For a batch,
the closure uses $W=QS/R$ for $R>0$, giving $\E[W]=pS$ and $S^+=S^--W$.
This preserves the conditional mean loss of $S$. Its higher joint moments
are an approximation. Feasibility repairs enforcing $R^2/M\leq S\leq R$
and bounded-variance limiting are recorded as numerical diagnostics.

For fixed interval counts, the first two moments of cumulative charge close
in $\E[R]$, $\E[S]$, $\E[R^2]$, $\E[Q_{\rm tot}]$,
$\E[Q_{\rm tot}R]$ and $\E[Q_{\rm tot}^2]$. Indeed, the charge mean is
linear in $R$, its second moment is linear in $S$ and $R^2$, and recovery
is affine in the individual cell recoveries. Preserving the conditional
charge moments and $\E[W]=pS$ therefore preserves those aggregate moments
for the discretized process, provided no feasibility repair or variance
limiting modifies the update. The beta distribution approximates the shape
and higher moments of that process; it does not introduce an additional
approximation to these closed first and second moments.

For mixed optical and dark triggers, labels are exchangeable. Conditional
on the interval charge and squared-charge sum, random assignment of $k_s$
signal labels among $k=k_s+k_d$ triggers gives
\begin{align}
 \E[Q_s\mid Q,W]&=\alpha Q,\qquad \alpha=k_s/k,\\
 \Var(Q_s\mid Q,W)&=\frac{\alpha(1-\alpha)}{k-1}(kW-Q^2),\quad k>1.
\end{align}
The prototype substitutes the closed estimate of $W$, constrained to its
conditional bounds for this label allocation, and samples bounded $Q_s$;
$Q_d=Q-Q_s$. This preserves the total sampled charge and both population
charge bounds. For $k$ trigger amplitudes in $[0,1]$ summing to $Q$, these
conditional bounds are $Q^2/k\leq W\leq\lfloor Q\rfloor+
(Q-\lfloor Q\rfloor)^2$. Optical/dark partition moments inherit this
approximation to $W$; the saved production signal uses their sum $Q$.

\paragraph{Arrival-time discretization.}
Dense buckets are split into $K$ equal intervals. Successive binomial
partitions give the exact multinomial allocation of the fixed optical and
dark counts to these intervals. Half an interval of exponential recovery
precedes the occupancy draw and half follows it. Collapsing arrivals within
an interval to its midpoint is the time-discretization approximation;
comparison of $K=32,64,128$ quantifies it. Buckets with at most 16 triggers
instead sample their individual uniform arrival times and apply successive
single-trigger moment draws. They retain the same compact state.

The TOML selector \code{response_model = "statistical"} chooses this path;
\code{statistical_substeps} and \code{sparse_trigger_limit} specify its
resolution. The prototype card uses 128 intervals per 250~ns bucket and a
sparse limit of 16 triggers. The microscopic path remains the reference. Model cards and code
hashes distinguish the saved responses. The validation script
\code{scripts/production/benchmark_sipm_statistical.jl} compares transient
means, variances, adjacent-bucket covariance, total event charge and runtime
for synthetic histories and the saved photons from both operating points.

\subsection{Inverse correction and what it recovers}
Solving equation~\eqref{eq:forward-mean} for $n$ gives AXEL's correction
\begin{equation}
 \widehat n=F(m)=\frac{m}{1-km},\qquad km<1.
 \label{eq:inverse}
\end{equation}
AXEL measured device recovery with an LED and a low-intensity reference,
then used individual recovery times in the nonlinearity correction. The
2024 study reports a mean of 73.4~ns for its devices, with 3600 microcells
and 200~ns sampling~\cite{Yoshida2024AXELResponse}. It also adjusts the
effective recovery calibration using photopeaks and a corrected squared-sum
observable; the 2025 analysis develops this fine-tuning alongside EL-gain,
time and drift-distance corrections~\cite{Akiyama:2025dia}. Our initial
forward-model $\tau$ is the adopted mean; a later reconstruction calibration
may fit an effective value to the simulated pulse population.

The correction maps measured charge to an estimate of the primary avalanche
count under its rate assumptions. It cannot identify which photons were
missed, recover the random microcell collision history, or remove its event
fluctuations. Even with exact mean calibration,
$\E[F(m_{\mathrm{event}})]$ need not equal $F(\E[m_{\mathrm{event}}])$.
For small fluctuations around $m_0$, its local noise amplification is
\begin{equation}
 F'(m_0)=\frac{1}{(1-km_0)^2},\qquad
 \Var(\widehat n)\simeq\frac{\Var(m)}{(1-km_0)^4}.
 \label{eq:inverse-noise}
\end{equation}
The positive curvature produces a leading mean shift
$k\Var(m)/(1-km_0)^3$ in this small-fluctuation approximation. Both
expressions fail near the pole, where the full response distribution must be
used. Correcting a clipped output cannot determine the lost current. Those
samples retain their saturation flag and require an explicitly justified
censored-data or event-selection treatment in a later analysis.

\subsection{Dark subtraction, PDE and energy analysis}
With stationary signal plus dark counts, equation~\eqref{eq:inverse}
estimates their \emph{total} primary count. In that limit a possible signal
estimator is
\begin{equation}
 \widehat n_s=F(m)-R_d\Delta t.
 \label{eq:dark-subtraction}
\end{equation}
Subtracting a full-charge mean dark signal before the nonlinear inversion is
generally inequivalent, because signal and dark avalanches compete for the
same recovering cells. For a time-varying event, dark-only calibration and
joint forward simulations should establish the subtraction and inversion
bias. Subtracting a known mean dark contribution leaves its statistical
variance. The saved signal/dark charge components permit direct checks of
these effects in simulation, while a real acquisition measures their sum.

After validating the correction, dividing $\widehat n_s$ by the calibrated
PDE estimates incident photons. The division changes the scale and preserves
the detection fluctuation. Converting incident light to EL electrons and
deposited energy adds the separate optical collection, EL gain and
attachment calibrations already associated with the upstream gas model.
These stages should be tested against saved truth quantities in order:
sensor charge, primary detected count, incident photons and event energy.
The inverse model is a subsequent analysis stage; it does not alter the
saved forward sensor or ASIC histories.

\section{Implementation, persisted stages and validation}
The photon reader is \code{read_photon_history} in
\code{src/event_analysis/production_readout.jl}, exported through
\code{EventAnalysis.jl}. It returns owned arrays with signed SiPM and bucket
indices. The \code{Sensor} type in
\code{src/ion_signal/sipm_response.jl} stores PDE, microcell count, gain,
recovery time and dark rate. Its \code{simulate_channel} function samples
successive photon buckets and returns owned count and charge arrays.
The \code{Electronics} type in \code{src/ion_signal/sipm_electronics.jl}
stores sampling interval, noise and range; \code{apply_electronics} returns
owned noise/current/flag arrays while preserving its input charge.
The diagnostic \code{axel_mean} and \code{axel_inverse} functions evaluate
the stationary formulas; the inverse reports a validity flag and rejects
clipped samples or a nonpositive denominator.

The sensor and electronics cards are independently versioned at
\code{config/specs/response/sipm/sipm_v1.toml} and
\code{config/specs/response/sipm_electronics/readout_v1.toml}.
The acquisition adapter and persisted stages are implemented in
\code{scripts/production/sipm_pipeline.jl}, with
\code{scripts/production/run_sipm_pilot.jl} as the Julia pilot entry point.
A production run card selects the inputs and stages. Canonical card hashing and event-keyed seeds follow
\code{scripts/production/specs.jl}. The sensor file schema is \code{nausicaa.sipm_file.v1}, with event schema
\code{nausicaa.sipm_history.v1}; the ASIC file uses
\code{nausicaa.electron_response_file.v1}, with event schema
\code{nausicaa.electron_response_history.v1}. Metadata identify the model,
algorithm, units and initialization convention explicitly.

For each saved channel--bucket cell, the sensor stage records incident
photons, primary signal and dark avalanche counts, signal and dark charge,
and total charge. The saved $q_0$ and $\Delta t$ determine PE-equivalent
response and bucket-average current exactly from the charge arrays. The output
must also describe the acquisition channel set and window so that absent
photon input cells can be represented with dark and electronics noise.
The electronics stage references the saved sensor identity and stores its
noise draw, current before and after clipping, and saturation flag. This
separation permits ASIC changes while preserving the identical sampled
sensor realization.

Source-content hashes, event and gas identifiers, parent-pair checksum,
geometry/time conventions, exact TOML text and canonical hashes, RNG seeds,
and code identities accompany the outputs. Output files are marked complete
only after successful writing. Physical calculations use Julia and save
results to HDF5; Python presentation code reads those results to generate
parameter tables and this report's numerical additions. The four source
files are preserved, and the response is first evaluated on a small pilot.

The validation program tests physical limits and persisted behavior:
\begin{enumerate}
\item Check the binomial mean and variance at negligible occupancy, and
one- and two-avalanche charge against equation~\eqref{eq:recovery}.
\item Check recovery through empty buckets and a signal at a bucket
boundary; changing bucket labels must not reset the microcells.
\item Check Poisson dark-count mean, variance and area scaling, with
charge approaching equation~\eqref{eq:dark-current} in the dilute limit.
\item Compare a constant-rate simulation after equilibration with the
stationary mean in equation~\eqref{eq:forward-mean}; then quantify deviations
for finite pulses and changing illumination.
\item Check equation~\eqref{eq:charge-preservation} and verify that saved
signal and dark components sum to the sensor response.
\item Check the Gaussian noise distribution, independence prescription,
upper clipping and retained negative samples, including zero illumination.
\item Round-trip both HDF5 stages, verify identities and event seeds, and
confirm that changing only ASIC parameters preserves the sensor realization.
\item On the four-input pilot, report channel/bucket occupancy, charge loss
from recovery, dark contribution, clipped fraction and resource use before
scaling to all events. Later inversion tests report bias and resolution as
functions of intensity, pulse shape and distance from clipping.
\end{enumerate}

The active approximations are uniform sub-bucket timing and illumination,
fixed PDE during recovery, common microcell gain and recovery time,
independent primary dark pulses, zero correlated avalanches, and a charge
bucket ASIC with independent Gaussian noise. They define a reproducible
baseline. Sensor measurements and ASIC information can replace each input
through the versioned cards while keeping the saved stage boundaries and
truth comparisons explicit.

\InputIfFileExists{latex_generated/sipm_response/pilot_results.tex}{}{}

The pilot demonstrates substantial clipping under the adopted 1~mA
bucket-average-current limit. The saved sensor stage retains the charge that
precedes this loss, allowing the current range and reconstruction strategy to
be reviewed independently. A recovery correction alone cannot restore the
clipped samples. The stationary and transient tests also show why the inverse
requires calibration for the time structure of the actual light pulses.
The completed 8,000-event bb0nu response and the operating-point comparison
are documented in \code{docs/response_overview.md}.

The statistical implementation uses the same production PE-per-bucket schemas:
\code{nausicaa.sipm_pe_file.v1} for sensor charge and
\code{nausicaa.sipm_readout_pe_file.v1} for ASIC response. Its card is
\code{config/specs/response/sipm/sipm_statistical_v1.toml}. The source,
interfaces, validation commands and two-event integration pilot are described
in \code{docs/sipm_statistical_implementation.md}. The generated comparison is
\code{latex_generated/sipm_statistical_validation/validation.pdf}, with an
HTML counterpart at \code{results/sipm_statistical_validation/validation.html}.

\bibliographystyle{unsrt}
\bibliography{notes/physics/references}
\end{document}
